\chapter{Introduction}

\section{Contexte}
Dans le cadre de notre projet P2 SP, nous avons pour but de réaliser un logiciel de traitement de flux vidéo en temps réel.

\section{Principe de notre logiciel}

Notre logiciel permettra de lire un flux vidéo en entrée, de le traiter en temps réel et de visualiser le résulat en même temps.
Nous pourrons notamment augmenter la saturation de l'image, changer les couleurs et appliquer divers filtres.
Notre logiciel s'adresse aux créateurs de contenu vidéo en direct, notamment les streameurs.

\begin{figure}[h]
  \centering
  \includegraphics[width=\textwidth]{./images/image_processing_exemple.jpg}
  \caption{Exemple de traitement de flux}
  \label{image_processing_exemple}
\end{figure}

Extrait des conseils pour écrire un rapport de P3 et de Bachelor présents \href{https://gitlab-etu.ing.he-arc.ch/isc/documentation/projet-p3-et-bachelor/-/wikis/help-report}{ici} :
\begin{quote}    
  \begin{itemize}   
    \item exposer le sujet (ou le thème, la matière prise en considération) en faisant valoir son importance et son originalité,
    \item articuler la problématique soulevée,
    annoncer le plan.
  \end{itemize}
\end{quote}

Pour inclure une image, il faut faire comme ceci :
\begin{figure}[h]
  \centering
  \includegraphics[width=\textwidth]{./images/logos/he-arc-banner.jpg}
  \caption{Teaser figure.}
  \label{Fig_teaser}
\end{figure}

On peut ensuite facilement référencer la Figure \ref{Fig_teaser} ailleurs dans le texte.

Pour rajouter une nouvelle figure et la référencer (voir Figure \ref{Fig_exemple}).
\begin{figure}[h]
    \centering
    \includegraphics[width=0.5\textwidth]{./images/logos/he-arc-logo.png}
    \caption{Ceci est le nouveau logo de la HE-Arc.}
    \label{Fig_exemple}
\end{figure}

À noter que la numérotation se fait automatiquement.
Les liens entre les références sont aussi gérés de manière transparente.

\section{Comment formater le texte}
\textbf{Comment mettre un texte en gras.}
\underline{Comment souligner un texte.} 
\textit{Comment mettre un texte en italique.}
\underline{\textbf{\textit{On peut aussi mélanger les 3.}}}

Pour commencer un nouveau paragraphe, il suffit d'ajouter 1 ligne vide le fichier .tex.

Pour faire une liste :
\begin{itemize}
    \item élément 1
    \item élément 2
    \item élément 3
\end{itemize}

Pour faire une liste numérotée :
\begin{enumerate}
    \item élément 1
    \item élément 2
    \item élément 3
\end{enumerate}

\section{Lorem ipsum}
\lipsum[1-1] % PLACEHOLDER
\subsection{Dolor sit}
\lipsum[1-1] % PLACEHOLDER
\subsubsection{amet}
\lipsum[1-1] % PLACEHOLDER
\paragraph{consectetuer}
\lipsum[1-1] % PLACEHOLDER